\chapter*{Özet}
\addcontentsline{toc}{chapter}{Özet}

%% Edit below this line
İşaret dili, milyonlarca işitme engelli birey için temel iletişim aracıdır;
ancak günlük hayatta tercüman bulunmaması, işiten dünya ile aralarında sürekli
bir iletişim engeli oluşturmaktadır. Bu proje, Amerikan İşaret Dili (ASL)
hareketlerini tanıyıp gerçek zamanlı olarak sesli çıktıya dönüştüren, kamera
veya herhangi bir dış altyapı gerektirmeyen, kendi kendine yeten giyilebilir
bir eldiven sunmaktadır.

Sistem, beş parmağın her birinin tırnağına kompakt bir eylemsizlik ölçüm
birimi (MPU6050) ve bir mikrodenetleyici (Seeed Studio XIAO ESP32-C3)
yerleştirir. Dört ``uydu'' modül, füzyonlanmış yönelim ve açısal hız
verilerini ESP-NOW protokolü üzerinden orta parmaktaki beşinci ``ana'' modüle
iletir. Ana modül, 25\,Hz hızında 45 kanallı bir öznitelik vektörü oluşturarak
bunu Bluetooth Düşük Enerji (BLE) üzerinden bir akıllı telefona aktarır. Bir
Flutter uygulaması, istek üzerine sabit uzunlukta bir sensör penceresi
yakalar, nicelenmiş tek boyutlu bir evrişimsel sinir ağını tamamen cihaz
üzerinde çalıştırır ve tanınan harf veya kelimeyi metinden konuşmaya çevirme
ile seslendirir.

Çalışmanın temel bulgularından biri, öznitelik temsilinin belirleyici
önemidir. Mutlak parmak yönelimine dayalı statik el şekli tanıma, eylemsizlik
referans çerçevesinin oturumlar arasında tam olarak yeniden üretilememesi ve
sensörlerin sınırlı açısal çözünürlüğünün yalnızca küçük başparmak
hareketleriyle ayrılan el şekillerini ayırt edememesi nedeniyle kırılgan
kalmıştır. Bu nedenle proje, ayırt edici bilgiyi kalibrasyondan bağımsız
açısal hız kanallarına taşıyan, harfe özgü belirgin \emph{hareketler} ile
\emph{mutlak} öznitelikleri benimsemiştir. Dağıtılan modeller 137\,KB flash ve
22.7\,KB RAM kullanır ve çıkarım başına yaklaşık 1\,ms sürede çalışır; kablosuz
bağlantı beş parmağın tüm akışlarını 25\,Hz hızında ihmal edilebilir paket
kaybıyla iletir. Proje, düşük maliyetli, yaygın donanım üzerinde doğru, düşük
gecikmeli ve tamamen cihaz içi işaret dili çevirisinin uygulanabilir olduğunu
göstermekte ve böyle bir sistemi yöneten mühendislik ödünleşimlerini—radyo
birlikte çalışabilirliği, güç bütünlüğü, termal davranış ve öznitelik
değişmezliği—belgelemektedir.

%% Until here
\vfill
%% Edit after {Anahtar Kelimeler:}
\textbf{Anahtar Kelimeler:} işaret dili tanıma, giyilebilir teknoloji,
eylemsizlik ölçüm birimi, ESP-NOW, Bluetooth Düşük Enerji, uç birim makine
öğrenmesi, tek boyutlu evrişimsel sinir ağı, hareket tanıma.
\clearpage
